\section{Related Work}

\textbf{Visual and point-based policies.}
ACT~\cite{zhao2023act} predicts action chunks, DP~\cite{chi2023dp} generates actions through conditional denoising, and $\pi_{0.5}$~\cite{physicalintelligence2025pi05} uses large-scale vision-language-action pretraining.
Geometric alternatives include DP3's sparse scene point clouds~\cite{ze2024dp3}, PointPolicy's semantic observation and action keypoints learned from human videos~\cite{haldar2025pp}, and PointBridge's unified points for synthetic-data training and sim-to-real transfer~\cite{haldar2026pb}.
GraphPoint organizes task-relevant points by semantic role and conditions their interaction on action types and modifiers.

\textbf{Object-centric and graph representations.}
GROOT~\cite{zhu2023groot} uses segmented object point clouds and segmentation correspondence for visual generalization.
GraphMimic~\cite{chen2025graphmimic} predicts object and visual-action graphs from videos to condition control.
Compose by Focus~\cite{qi2026compose} encodes task-relevant subgraphs and combines diffusion-based atomic skills with a vision-language planner.
Building on these established object-centric and graph-based approaches, GraphPoint couples actor/patient/target roles with explicit action/modifier conditioning and point-trajectory generation.

\textbf{Compositional generalization.}
CLIPort~\cite{shridhar2022cliport} combines semantic features with spatial action prediction; VIMA~\cite{jiang2023vima} evaluates multimodal prompting across four generalization levels.
CompoSuite~\cite{mendez2022composuite} holds out combinations of robots, objects, obstacles, and objectives, while LIBERO~\cite{liu2023libero} studies declarative and procedural transfer and CALVIN~\cite{mees2022calvin} evaluates language-conditioned task sequences.
CoMani isolates modifier composition, action-type composition, and three-step execution in LIBERO, using matched scenes to reduce visual task shortcuts and atomic-only training to test sequential reuse.
